\chapter{Marco Teórico}

La noción de \textit{typestate} fue introducida en 1986 por Strom y Yemini~\cite{typestate}. A pesar de ya en esa época existir mecanismos de prevención de errores en programas, tales como los sistemas de tipado o reglas basadas en scopes estáticos, los autores identificaron una limitación en los mismos. Específicamente, fallaban en detectar varias acciones potencialmente prohibidas en ciertos contextos, tales como desreferenciar un puntero apuntando a memoria ya liberada, o usar una variable antes de ser inicializada. Mientras que los tipos restringen las operaciones permitidas para un dato/objeto, y los scopes el acceso a variables fuera de su ciclo de vida, los typestates intentan restringir el orden en el cual las operaciones sucedían bajo un contexto particular.

El funcionamiento de los mismos es de la siguiente manera. Dado un tipo, se le asocia un conjunto de typestates, y se define para cada uno un conjunto de operaciones permitidas. A su vez, se definen reglas de transición entre ellos a través de estas operaciones. En cualquier punto del programa, un objeto estará en alguno de estos typestates y podrán ser aplicadas solamente sus operaciones permitidas. Al hacerlo, el objeto potencialmente transicionará a otro typestate según las reglas definidas, cambiando el conjunto de operaciones disponibles.

El paper menciona además conceptos que hoy en día \textit{JaTyC} implementa, tales como el de \textbf{completitud de protocolo}, donde se espera que los objetos lleguen a un estado final antes de que el programa termine. Según los autores, una ejecución es \textit{typestate correcta} si y solo si todas las operaciones realizadas en objetos son permitidas por su typestate, y todos los objetos completan su protocolo  .

In the paper, the notion of protocol completion also makes an initial appearance, where objects are supposed to return to the initial typestate before the program terminates for it to have a *typestate-correct* execution.

The formal theory for typestates and an initial algorithm to determine statically if a program is *typestate-consistent* is also presented in this paper.

% - Talk bout the paper `Lightweight Object Specification with Typestates` which uses typestates to define object protocols even further.
% - Perhaps we could add more papers here if necessary.